\chapter{한 위치의 판정: \texorpdfstring{\code{bombe\_core}}{bombe\_core}}\label{ch:rtl-core}

\code{bombe\_core}는 드럼 위치 하나에 대해 봄베가 하는 일 전부를 한다. 메뉴대로 스크램블러를 걸고, 전류를
퍼뜨리고, 테스트 레지스터를 읽어 정지 여부를 알린다. Rust의 \code{propagate}에 해당한다.

\section{가장 큰 설계 결정: 고리를 레지스터로 끊는다}

실제 봄베의 배선망은 전선이다. 전압을 걸면 전류는 이어진 선 전체로 사실상 즉시 퍼지고, 드럼이 움직여 접점이
떨어지면 사라진다. 이것을 게이트로 그대로 옮기려는 것이 첫 번째 유혹이다.

\Why 그대로 옮길 수 없는 이유는 배선이 \emph{양방향}이기 때문이다. 스크램블러는 케이블 $a$의 선을 케이블 $b$의
선에 잇는다. 게이트로 쓰면 ``$b$의 선은 $a$의 선들의 OR''이고 동시에 ``$a$의 선은 $b$의 선들의 OR''이다.
출력이 자기 입력으로 돌아오는 고리, \term{조합 루프}{combinational loop}가 된다
(그림~\ref{fig:g-loop}). 이 고리에는 문제가 둘 있다.

\begin{itemize}
\item \textbf{꺼지지 않는다.} OR 게이트 두 개가 서로의 출력을 입력으로 받으면, 한번 1이 된 뒤에는 바깥 입력을
거두어도 서로를 1로 붙들고 있다. 전선의 전류는 접점이 떨어지면 사라지지만 게이트의 고리는 의도하지 않은
기억 소자가 된다. 다음 드럼 위치를 시험하려면 선을 모두 0으로 되돌려야 하는데 방법이 없다.
\item \textbf{시간을 따질 수 없다.} \ref{ch:sync}장의 규칙은 ``레지스터에서 레지스터까지의 지연이 클럭
주기보다 짧아야 한다''였다. 고리에는 출발 레지스터도 도착 레지스터도 없다. 도구가 지연을 계산할 수 없으므로
회로가 제 속도로 도는지 확인할 길이 없다.
\end{itemize}

\begin{figure}[ht]
\centering
\begin{tikzpicture}[x=1cm,y=1cm]
  \node[font=\small\bfseries] at (2.0,2.2) {조합 루프 (쓸 수 없다)};
  \node[or gate US,gate] (o1) at (1.2,1.0) {};
  \node[or gate US,gate] (o2) at (3.0,-0.4) {};
  \draw (o1.input 1) -- ++(-0.8,0) node[left,font=\footnotesize] {바깥 입력};
  \draw[cred,line width=0.9pt] (o1.output) -- ++(0.5,0) |- (o2.input 1);
  \draw (o2.input 2) -- ++(-1.2,0) node[left,font=\footnotesize] {0};
  \draw[cred,line width=0.9pt] (o2.output) -- ++(0.4,0) |- (1.6,-1.3) -| (0.2,0.5) |- (o1.input 2);
  \node[note,text=cred] at (2.0,-1.9) {한번 1이 되면 서로를 붙든다};
  \begin{scope}[xshift=7.6cm]
  \node[font=\small\bfseries] at (2.9,2.2) {레지스터로 끊은 구조};
  \node[reg,minimum height=15mm,minimum width=15mm] (r) at (0,0.2) {\code{live\_q}};
  \pic at (r.south) {clkmark};
  \node[cloud,minimum width=36mm,minimum height=15mm] (f) at (3.9,0.2) {한 걸음 퍼뜨리기};
  \draw[->] (r) -- (f);
  \draw[->] (f.east) -- ++(0.5,0) |- (2.0,-1.3) node[pos=0.75,below,font=\footnotesize] {\code{live\_n}} -| ([xshift=-7mm]r.west) -- (r.west);
  \node[note] at (2.9,-2.2) {지우려면 레지스터에 0을 넣으면 된다};
  \end{scope}
\end{tikzpicture}
\caption{고리와 그 해결.}
\label{fig:g-loop}
\end{figure}

해결은 고리의 한 곳을 레지스터로 끊는 것이다. 676가닥의 상태를 레지스터 \code{live\_q}에 담고, 조합 논리는
``지금 상태에서 \emph{한 걸음} 더 퍼진 상태'' \code{live\_n}만 계산한다. 클럭이 뛸 때마다 그 결과를 레지스터에
다시 담는다. 더 변하지 않으면 전류가 다 퍼진 것이다. 새 위치를 시작할 때는 레지스터에 0을 넣고 전압을 거는
선 하나만 1로 둔다.

\General 이런 꼴의 계산을 \term{고정점 반복}{fixed-point iteration}이라 한다. 같은 규칙을 되풀이 적용하다가
결과가 더 변하지 않으면 멈춘다. 소문이 퍼지는 것을 하루 단위로 보면 같다. 매일 아침, 어제까지 소문을 들은
사람 \emph{모두}가 아는 사람들에게 동시에 전한다. 어느 날 새로 들은 사람이 없으면 소문은 다 퍼진 것이다.

\Ours 이 구조가 실제 기계와 같은 답을 내는 까닭은 \ref{sec:propagate}절에서 보았다. 구하는 것은 전압을 건
선에서 닿을 수 있는 선의 집합이고, 그 집합은 살피는 순서와 무관하다. Rust는 한 가닥씩, 회로는 클럭마다
전부를 한 걸음씩 살피지만 도달하는 상태는 같다. 이 주장은 시험으로도 확인한다. 회로의 676비트를 Rust의
676비트와 통째로 비교한다.

\section{입출력}

\begin{lstlisting}[language=SV,caption={\code{bombe\_core.sv}의 입출력(로터 순서 입력은 \code{scrambler}와 같아 생략).}]
module bombe_core #(parameter int MAX_EDGES = 16) (
  input  logic                       clk, rst_n,
  ...
  // menu
  input  logic [MAX_EDGES-1:0]       edge_en,
  input  logic [MAX_EDGES-1:0][4:0]  edge_a,
  input  logic [MAX_EDGES-1:0][4:0]  edge_b,
  input  logic [MAX_EDGES-1:0][7:0]  edge_pos,
  input  logic [4:0]                 test_letter, input_wire,
  input  logic                       diagonal_en,
  input  logic                       turnover_en,
  input  logic [7:0]                 turnover_at,
  input  logic                       early_exit,
  // drum position (core offsets at the first crib letter)
  input  logic [4:0]                 pos_g, pos_l, pos_m, pos_r,
  // control
  input  logic                       start,
  output logic                       busy, done, stop,
  output logic [25:0]                test_register,
  output logic [26*26-1:0]           live,
  output logic [15:0]                iterations
);
\end{lstlisting}

메뉴가 어떻게 들어오는지를 보자. 소프트웨어에서 메뉴는 길이가 정해지지 않은 목록이었다. 회로에는 길이가
변하는 것이 없다. 그래서 선을 걸 수 있는 \emph{자리}를 \code{MAX\_EDGES}개(기본 16개) 만들어 두고, 자리마다
쓰는지 여부(\code{edge\_en}), 양 끝 글자(\code{edge\_a}, \code{edge\_b}), 크립 위치(\code{edge\_pos})를 준다.

\Why 회로의 부품 수는 만들 때 정해진다. ``선이 14개인 메뉴''와 ``선이 9개인 메뉴''를 같은 회로로 다루려면
가장 큰 경우에 맞춰 자리를 만들고 안 쓰는 자리를 끄는 수밖에 없다.

\Ours 실제 봄베가 꼭 이렇다. 스크램블러는 36개가 고정으로 달려 있고, 조작자는 메뉴에 필요한 만큼만
케이블로 이어 쓴다. \code{MAX\_EDGES}는 기계에 달린 스크램블러의 수이고 \code{edge\_en}은 그 스크램블러에
케이블이 꽂혀 있는가다.

제어 신호는 셋이다. \code{start}를 한 클럭 올리면 판정을 시작하고, 끝나면 \code{done}이 한 클럭 올라온다.
그때 \code{stop}이 1이면 정지다.

\section{스크램블러를 선마다 하나씩}

\begin{lstlisting}[language=SV,caption={\code{bombe\_core.sv}: 선마다 스크램블러 하나(일부).}]
for (genvar e = 0; e < MAX_EDGES; e++) begin : g_edge
  ...
  always_comb begin
    sum_r    = {4'd0, pos_r} + {1'b0, edge_pos[e]};
    mod_r    = sum_r % 9'd26;
    core_r_e = mod_r[4:0];
    core_m_e = (turnover_en && edge_pos[e] >= turnover_at)
               ? wirings_pkg::add26(pos_m, 5'd1) : pos_m;
  end
  scrambler u_scrambler (
    ...
    .core_g (pos_g), .core_l (pos_l), .core_m (core_m_e), .core_r (core_r_e),
    .perm   (perm_c[e*130 +: 130])
  );
end
\end{lstlisting}

\code{for (genvar ...)}는 \term{생성 반복문}{generate loop}이다. \ref{ch:hdl}장의 반복문과 마찬가지로 시간을
쓰지 않고 부품을 복사한다. 여기서는 \code{scrambler} 모듈을 16개 놓는다. \code{scrambler u\_scrambler (...)}는
부품 하나를 놓고 단자를 잇는 문장이다. \code{.core\_m(core\_m\_e)}는 ``부품의 \code{core\_m} 단자에 전선
\code{core\_m\_e}를 잇는다''는 뜻이다.

선마다 다른 것은 오른쪽 코어 위치뿐이다. 드럼 위치 \code{pos\_r}에 그 선의 크립 위치를 더한다. Rust의
$(r + k) \bmod 26$이다. 넘어감 가설이 켜져 있고 크립 위치가 가설 이후이면 가운데 코어 위치에 1을 더한다.

두 가지를 짚어 둔다. 첫째, \code{sum\_r \% 9'd26}에는 앞 장에서 피하라고 한 \code{\%}가 있다. 크립 위치는
0--255일 수 있어 ``26을 한 번 빼기''로는 안 되기 때문인데, 시뮬레이션용이라 그대로 두었다. FPGA 엔진에서는
이런 계산을 구동 전에 미리 해 둔다. 둘째, 이 구조는 \ref{sec:tradeoff}절에서 본 대로 ``위치마다 전부 다시
계산''이다. Rust 모델보다도 단순하다. 표를 미리 만들지 않고 16개의 스크램블러가 늘 계산하고 있다.

\section{전파 논리}

코드~\ref{lst:g-spread}의 전파 논리를 다시 읽는다. 이번에는 대각선 보드까지 본다.

\begin{lstlisting}[language=SV,caption={\code{bombe\_core.sv}: 한 걸음 퍼뜨리기.},label={lst:spread-full}]
always_comb begin
  live_n = live_q;
  y      = '0;
  for (int i = 0; i < MAX_EDGES; i++) begin
    cable_a = 26 * edge_a[i];
    cable_b = 26 * edge_b[i];
    if (edge_en[i]) begin
      for (int x = 0; x < 26; x++) begin
        y = {27'd0, perm_q[(i*26 + x)*5 +: 5]};
        if (live_q[cable_a + y]) live_n[cable_b + x] = 1'b1;
        if (live_q[cable_b + y]) live_n[cable_a + x] = 1'b1;
      end
    end
  end
  if (diagonal_en) begin
    for (int c = 0; c < 26; c++) begin
      for (int w = 0; w < 26; w++) begin
        if (live_q[c*26 + w]) live_n[w*26 + c] = 1'b1;
      end
    end
  end
end
\end{lstlisting}

\begin{itemize}
\item 2행. 다음 상태는 현재 상태에서 출발한다. 한 번 켜진 선은 꺼지지 않는다.
\item 4--14행. 메뉴의 선 \code{i}마다, 건너편 케이블의 선 \code{x}마다: 치환표에서 \code{y = perm[x]}를
읽는다. 이쪽 케이블의 \code{y}번 선이 켜져 있으면 건너편의 \code{x}번 선을 켠다. 스크램블러의 치환은 두 번
적용하면 원래대로 돌아오므로(반사판 때문에), ``$y$가 $x$로 간다''와 ``$x$가 $y$로 간다''가 같은 표로 표현된다.
그래서 10행과 11행이 같은 \code{y}를 쓴다. 두 줄은 전류의 두 방향이다.
\item 15--21행. 대각선 보드. 케이블 $c$의 $w$번 선이 켜져 있으면 케이블 $w$의 $c$번 선을 켠다. 여기서
\code{c}와 \code{w}는 반복 변수라 합성할 때 상수가 된다. 그러므로 이 부분은 선택 회로가 아니라 고정된 전선
연결과 OR이다. 실제 봄베의 대각선 보드가 \emph{영구 배선}이었던 것과 정확히 같다.
\end{itemize}

\Why 주의할 것은 읽는 쪽이 모두 \code{live\_q}(현재 상태)이고 쓰는 쪽이 모두 \code{live\_n}(다음 상태)이라는
점이다. 이번 걸음에 새로 켜진 선은 이번 걸음에서는 이웃을 켜지 못한다. 다음 클럭에 \code{live\_q}가 된 뒤에
켠다. 그래서 정확히 ``한 걸음''이다. 만약 \code{live\_n}을 읽었다면 한 클럭 안에서 여러 걸음이 이어져
앞에서 피하려던 고리가 되살아난다.

4--14행의 \code{live\_q[cable\_a + y]}가 \ref{ch:lesson}장의 그 코드다. 대각선 보드의
\code{live\_q[c*26 + w]}와 생긴 것은 비슷하지만 회로로는 전혀 다르다. 한쪽은 번호가 상수여서 전선 한 가닥이고,
다른 쪽은 번호가 실행 중에 정해져서 676:1 선택 회로다. 같은 꼴의 코드가 번호가 상수인가 변수인가에 따라
수백 배 다른 크기가 된다.

\section{제어: 상태 기계}

\begin{lstlisting}[language=SV,caption={\code{bombe\_core.sv}: 제어(리셋 부분 생략).},label={lst:core-fsm}]
always_ff @(posedge clk or negedge rst_n) begin
  ...
    done <= 1'b0;
    case (state)
      IDLE: begin
        if (start) begin
          perm_q     <= perm_c;
          live_q     <= '0;
          live_q[26*test_letter + input_wire] <= 1'b1;
          iterations <= '0;
          state      <= SPREAD;
        end
      end
      SPREAD: begin
        live_q     <= live_n;
        iterations <= iterations + 16'd1;
        if (live_n == live_q || (early_exit && test_n == FULL)) begin
          state <= IDLE;
          done  <= 1'b1;
        end
      end
    endcase
end

assign test_register = test_q;
assign stop          = (test_q != FULL);
\end{lstlisting}

\ref{sec:fsm}절에서 본 두 상태짜리 상태 기계다.

\begin{itemize}
\item 3행. \code{done}은 기본으로 0이다. 아래에서 1을 넣는 클럭에만 1이 된다. 같은 레지스터에 \code{<=}가 두
번 나오면 나중 것이 이긴다는 규칙을 이용해 ``한 클럭짜리 신호''를 만드는 흔한 방법이다.
\item 6--11행. \code{IDLE}에서 \code{start}를 보면: 스크램블러들이 내놓고 있는 치환표를 레지스터
\code{perm\_q}에 붙잡는다. 상태를 모두 0으로 하고(8행) 전압을 거는 선 하나만 1로 한다(9행). 8행과 9행도
같은 규칙이다. 나중에 나온 9행이 그 한 비트에 대해 이긴다.
\item 15--16행. \code{SPREAD}에서는 클럭마다 다음 상태를 받아들이고 걸음 수를 센다.
\item 17--20행. 끝내는 조건은 둘이다. 다음 상태가 현재 상태와 같거나(고정점), 조기 종료가 켜져 있고 테스트
레지스터가 다 찼거나.
\item 26행. 정지 판정. Rust의 \code{if reg != FULL}과 같은 한 줄이다.
\end{itemize}

그림~\ref{fig:g-core-timing}은 이 상태 기계가 한 위치를 판정하는 동안의 신호다.

\begin{figure}[ht]
\centering
\begin{tikztimingtable}[timing/slope=0.1,timing/coldist=4pt,xscale=2.6,yscale=1.3,semithick]
  \code{clk}      & 9{.5L .5H} \\
  \code{start}    & .5L 1H 7.5L \\
  \code{state}    & 1.5D{IDLE} 4D{SPREAD} 3.5D{IDLE} \\
  \code{live\_q}  & 1.5U 1D{$L_0$} 1D{$L_1$} 1D{$L_2$} 4.5D{$L_3$} \\
  \code{live\_n}  & 1.5U 1D{$L_1$} 1D{$L_2$} 5.5D{$L_3$} \\
  \code{iterations} & 1.5U 1D{0} 1D{1} 1D{2} 1D{3} 3.5D{4} \\
  \code{done}     & 5.5L 1H 2.5L \\
\extracode
  \begin{pgfonlayer}{background}
    \vertlines[cgray!50,dashed]{0.5,1.5,...,8.5}
  \end{pgfonlayer}
\end{tikztimingtable}
\caption{\code{bombe\_core}의 한 위치(개념도). $L_0$은 선 하나만 켜진 상태, $L_i$는 $i$걸음 퍼진 상태다.
네 번째 클럭에 \code{live\_n}과 \code{live\_q}가 같아져 다음 에지에서 \code{done}이 올라온다.}
\label{fig:g-core-timing}
\end{figure}

마지막 한 클럭은 ``더 변하지 않는다''는 것을 확인하는 데 쓰인다. 시뮬레이션에서 구동 중(조기 종료 사용)의
평균은 위치당 약 12클럭이었고, 끝까지 퍼뜨릴 때의 걸음 수는 13에서 25 사이였다.

\paragraph{리셋.} \code{negedge rst\_n}은 리셋 신호다. 전원이 들어온 직후 레지스터의 값은 정해져 있지 않으므로,
리셋으로 상태 기계를 \code{IDLE}에 놓는다. 프로그램에서 변수를 초기화하는 것에 해당하지만, 회로에서는 그
초기화도 전선(리셋 신호)을 통해 해야 한다.

\section{이 모듈의 시험}

테스트벤치는 Rust 코어가 내보낸 벡터 파일을 읽는다. 벡터 하나는 메뉴, 로터 순서, 드럼 위치, 그리고 그
위치에서의 676비트 상태와 정지 여부다. 벡터마다 두 번 돌린다. 조기 종료를 끄고 676비트 상태 \emph{전체}와
정지 여부를 비교하고, 조기 종료를 켜고 정지 여부와 테스트 레지스터를 비교한다. 벡터에는 넘어감 가설이 있는
경우, 대각선 보드를 끈 경우, M4의 경우가 들어 있다.
